\chapter*{Preliminaries and Notation}
\addcontentsline{toc}{chapter}{Preliminaries and Notation}

We assume the basic facts about sets, groups, rings, and modules, and about
categories, functors, and natural transformations.

\section{Sizes of sets and categories}\label{sec:P.1}

The natural numbers are $\mathbb{N}=\{0,1,2,\ldots\}$, and we write
$\mathbb{Z}$ and $\mathbb{Q}$ for the ring of integers and the field of
rational numbers, respectively. A family indexed by a set $I$ is written
$(X_i)_{i\in I}$. A finite family is a family whose index set $I$ is finite.
Finiteness of each $X_i$ is a condition separate from finiteness of the index
set.

Sizes of sets are handled in classical set theory with the axiom of choice,
together with Grothendieck universes. We fix universes
$\mathbb{U}\in\mathbb{V}$ as needed, and call a set that belongs to
$\mathbb{U}$ a $\mathbb{U}$-small set. We write
${\mathbf{Set}}_{\mathbb{U}}$ for the category of small sets and maps between
them, and treat this category itself in the larger universe $\mathbb{V}$. Where the
universe in use is clear, we omit the subscript and write ${\mathbf{Set}}$.

A category $\mathcal{C}$ is $\mathbb{U}$-small if its objects and its
morphisms form $\mathbb{U}$-small sets. The category $\mathcal{C}$ is locally $\mathbb{U}$-small if the set of
morphisms between each pair of objects $X,Y$ is $\mathbb{U}$-small, and
essentially $\mathbb{U}$-small if it is equivalent to a $\mathbb{U}$-small
category. Constructions that involve functor
categories use a universe large enough to contain their objects and
morphisms.

\section{Maps, morphisms, and composition}\label{sec:P.2}

The composite of maps $f:X\to Y$ and $g:Y\to Z$ is written
\[
gf=g\circ f:X\longrightarrow Z,
\qquad (gf)(x)=g(f(x)).
\]
The map on the right acts first. We use the same order of composition for
morphisms of categories and for functors. The identity morphism of an object
$X$ is written $\mathrm{id}_X$, and the identity functor of a category
$\mathcal{C}$ is written $\mathrm{Id}_{\mathcal{C}}$.

For a map $f:X\to Y$ and subsets $S\subseteq X$ and $T\subseteq Y$, we write
$f(S)$ for the image, $f^{-1}(T)$ for the preimage, and $f|_S:S\to Y$ for the
restriction. The fiber over $y\in Y$ is
\[
X_y=f^{-1}(\{y\})=\{x\in X\mid f(x)=y\}.
\]
A bijection is written $f:X\xrightarrow{\sim}Y$, and its inverse map is
written $f^{-1}:Y\to X$.

For a category $\mathcal{C}$, we write $\mathrm{Ob}(\mathcal{C})$ for its
collection of objects and $\mathrm{Hom}_{\mathcal{C}}(X,Y)$ for the set of
morphisms from $X$ to $Y$. When the category is clear, the latter is
abbreviated to $\mathrm{Hom}(X,Y)$. A morphism $f:X\to Y$ is an isomorphism if
there is a morphism $g:Y\to X$ with $gf=\mathrm{id}_X$ and
$fg=\mathrm{id}_Y$. Such a $g$ is unique, and we denote it by $f^{-1}$. We write
$X\cong Y$ for isomorphic objects, and $\mathrm{Aut}_{\mathcal{C}}(X)$ for the
automorphism group of $X$.

Commutativity of a diagram means that the composites along any two paths with
the same domain and codomain are equal. For example, the commutativity
condition for the square formed by $f:X\to Y$, $f':X'\to Y'$,
$a:X\to X'$, and $b:Y\to Y'$ is $bf=f'a$. When a comparison is made through a natural isomorphism, that natural
isomorphism is specified as part of the data of the diagram.

\section[Functors, natural transformations, and equivalences of
categories]{Functors, natural transformations,\texorpdfstring{\\}{ }and
equivalences of categories}\label{sec:P.3}

The actions of a functor $F:\mathcal{C}\to\mathcal{D}$ on objects and on
morphisms are written $F(X)$ and $F(f)$, respectively. The opposite category is written
$\mathcal{C}^{\mathrm{op}}$, and a contravariant functor is regarded as a
functor whose domain is $\mathcal{C}^{\mathrm{op}}$.

A natural transformation $\alpha:F\Rightarrow G$ between functors
$F,G:\mathcal{C}\to\mathcal{D}$ is a family that assigns to each object $X$ a
morphism $\alpha_X:F(X)\to G(X)$ such that, for each morphism $f:X\to Y$,
\[
G(f)\alpha_X=\alpha_YF(f).
\]
If every $\alpha_X$ is an isomorphism, $\alpha$ is called a natural
isomorphism and is written $\alpha:F\xRightarrow{\sim}G$. The composite of
natural transformations $\alpha:F\Rightarrow G$ and $\beta:G\Rightarrow H$ is
given by $(\beta\alpha)_X=\beta_X\alpha_X$.

A functor $F:\mathcal{C}\to\mathcal{D}$ is fully faithful if, for every pair
of objects $X,Y$, the map
\[
F_{X,Y}:\mathrm{Hom}_{\mathcal{C}}(X,Y)
\longrightarrow\mathrm{Hom}_{\mathcal{D}}(F(X),F(Y)),
\qquad f\longmapsto F(f)
\]
is a bijection. If every $D\in\mathrm{Ob}(\mathcal{D})$ is isomorphic to some
$F(X)$, $F$ is called essentially surjective.

An equivalence of categories $\mathcal{C}\simeq\mathcal{D}$ is given by
functors $F:\mathcal{C}\to\mathcal{D}$ and $G:\mathcal{D}\to\mathcal{C}$
together with natural isomorphisms
\[
\eta:\mathrm{Id}_{\mathcal{C}}\xRightarrow{\sim}GF,
\qquad
\varepsilon:FG\xRightarrow{\sim}\mathrm{Id}_{\mathcal{D}}.
\]
We call $G$ a quasi-inverse of $F$. For the standard definitions of
categories, functors, natural transformations, and equivalences of categories
we follow
\cite[\href{https://stacks.math.columbia.edu/tag/0013}{Tag~0013}]{Stacks}.

\section{Notation for coefficients and algebra}\label{sec:P.4}

A coefficient ring $k$ is a commutative ring with identity, and ring
homomorphisms preserve the identity. We write ${\mathbf{CommRing}}$ for the
category of commutative rings and ${\mathbf{Ab}}$ for the category of abelian
groups. Modules over $k$ are unital, $1m=m$, and we write $\mathrm{Mod}_k$ for
their category. A commutative $k$-algebra is a pair consisting of a commutative ring $A$
and a structure homomorphism $k\to A$. We allow the zero ring. These
conventions for commutative rings with identity and unital modules agree with
\cite[\href{https://stacks.math.columbia.edu/tag/0006}{Tag~0006},
\href{https://stacks.math.columbia.edu/tag/00AQ}{Tag~00AQ}]{Stacks}.

We write $A/I$ for the quotient ring of a ring $A$ by an ideal $I$, and $M/N$
for the quotient module of a module $M$ by a submodule $N$. The set of linear
maps between modules is written $\mathrm{Hom}_k(M,N)$, and the kernel and the
image of a homomorphism $u$ are written $\ker u$ and $\mathrm{im}\,u$. Abelian
groups and modules are written additively, and the zero element and the zero
map are both written $0$, according to context. General groups and
automorphism groups are written multiplicatively, with identity element $1$.

The tensor product over a fixed coefficient ring is written $M\otimes_k N$. A
change of coefficient ring is specified by a ring homomorphism
$\varphi:k\to k'$, and the extension of scalars of a module along it is
written $k'\otimes_k M$.

\section{Numbering and references}\label{sec:P.5}

Definitions, constructions, theorems, lemmas, propositions, corollaries, and
examples are numbered in one common sequence within each chapter of the main
body. For example, ``Definition~2.1'', ``Proposition~2.2'', and
``Definition~2.3'' all refer to items of Chapter~\ref{chap:2}. Sections, equations,
figures, and tables are each numbered within each chapter, and are referred to
as ``\S2.1'', ``(2.1)'', ``Figure~2.1'', and ``Table~2.1''. In the
Preliminaries the letter P, and in the appendices the letters A and B, are
used in place of the chapter number.

A citation gives a key such as \cite{Stacks}
together with the section or theorem number cited. Citations of the Stacks
project also give, in addition to section numbers, its permanent identifiers
called Tags. The bibliographic data are collected in the References.
